\documentclass{article}
\usepackage[T1]{fontenc}
\usepackage{lmodern}
\usepackage{graphicx}
\usepackage{amsmath,amssymb}
\usepackage[a4paper,margin=2cm]{geometry}
\usepackage[absolute,overlay]{textpos}
\setlength{\TPHorizModule}{1mm}
\setlength{\TPVertModule}{1mm}
\usepackage[indonesian]{babel}
\usepackage{hyperref}
\setlength{\emergencystretch}{2em}
% unit: Halaman 86

\begin{document}

% === Halaman 86 (python/native) ===

86

C. Affine Cipher

•

Merupakan perluasan Caesar cipher


Enkripsi: C  mP + b (mod n) 


Dekripsi: P   m–1 (C – b) (mod n)

•

Kunci: m dan b  

Keterangan:

1.

n adalah ukuran alfabet 

2.

m bilangan bulat yang relatif prima dengan n 

3.

b adalah jumlah pergeseran 

4.

Caesar cipher adalah khusus dari affine cipher dengan  m = 1 

5.

m–1 adalah inversi m (mod n), yaitu m  m–1  1 (mod n)

\end{document}
